\section{Introduction to the Environment}

\begin{frame}{What Will We Learn?}
	\begin{itemize}
		\item What \textbf{\texttt{self.env}} is and why every Odoo method uses it.
		\item The main parts of the environment: user, company, context, cursor, and registry.
		\item How to switch environment safely with \texttt{with\_context}, \texttt{with\_user}, \texttt{sudo()}, and friends.
		\item How to resolve XML IDs with \texttt{env.ref()}.
		\item The most important \textbf{\texttt{@api}} decorators and when to use each one.
	\end{itemize}
\end{frame}

\begin{frame}{Why the Environment Matters}
	\begin{itemize}
		\item In Odoo, business code almost never talks to PostgreSQL directly.
		\item Every recordset carries an \textbf{environment}.
		\item The environment answers questions such as:
		\begin{itemize}
			\item Who is the current user?
			\item Which company is active?
			\item What language / context flags are active?
			\item Which database cursor should be used?
		\end{itemize}
		\item Understanding \texttt{self.env} is the foundation for correct ORM code.
	\end{itemize}
\end{frame}
